% Fragmento público generado desde propietarios íntegros.
% Residencia: 52.11.2.
% Fuente: ../incoming/radion_monografia_20260827/manuscrito/sections/03_cierre_exacto.tex:286-309
\subsubsection{Caracterización estructural de \(\pi\) por compatibilidad de secciones}

Reordenar \eqref{eq:cierre-unitario} da
\begin{equation}
\boxed{
2\pih=
\frac{1+\Phi_T-\epsone}
{\frac5{36}+\frac{25}{9}\alphah}.}
\label{eq:pi-radion}
\end{equation}
Esta igualdad puede llamarse una nueva definición de \(\pi\) dentro de la
carta del radión si se entiende \emph{definición estructural} como una
caracterización exacta por compatibilidad de secciones. No es un segundo
generador independiente: \(\Delta_4\) y el operador de acarreo ya reciben la
sección coinductiva de \(\pi\). La novedad está en que la vuelta queda
caracterizada por el empalme entre publicación directa, torsión y memoria.

\begin{narrativebox}
La ecuación \eqref{eq:pi-radion} no dice que primero desconozcamos \(\pi\) y
lo adivinemos cerrando un ángulo. Dice que el \(\pi\) generado por la conexión
nonádica es exactamente el que hace compatibles dos lecturas internas del
mismo estado. Es una identidad de reconocimiento intrínseco, no una
calibración metrológica.
\end{narrativebox}
